\section{The Operator Is Already Traceable}\label{sec:turn-chain}

%% Numbers: empirics-futon/gen_turn_chain.py emits BOTH this section's macros
%% and the figure, from one live join, so prose and figure cannot drift (WR-8).

The closure argued in \S\ref{sec:operator} needs the operator to be
\emph{observable} before it can be modelled. That is an empirical claim about
this stack, not a design intention, so it is established separately and first.

Agent-authored commits record system output; they do not by themselves
measure operator attention. What records the operator is the evidence landscape, and it records
him at the granularity of the \emph{turn}.

Three mechanisms, built for unrelated reasons, bear on the same question.

\noindent\textbf{Route A: clock-in.} A turn taken while a mission is clocked carries
that mission on the turn record. This is not fully reliable --- the clock is
sometimes not set, and the gap is visible rather than hidden.

\noindent\textbf{Route B: retrieval.} The retrieval route embeds turns and matches them against
the pattern library, and the result is deposited as a \texttt{context-retrieval}
evidence event carrying pattern ids and cosine scores. Each recorded event carries \ChainCandidates{} ranked candidates, but the
operator's interface commits to the top one --- it resolves that pattern's sigil
and shows it beside the turn. The counts and the figure below therefore use rank
\ChainTopK{} only. Counting all \ChainCandidates{} would inflate every figure in
this section by roughly threefold while describing something no one ever saw.

\noindent\textbf{Route C: mention.} A mission document may simply name a
pattern in its text. This needs no turn and no clock at all; it is a property of
the corpus, and \ChainMissionDocs{} mission documents were scanned for it.

Each route was built independently; none was designed to supply the relation
measured by either of the others. Route~A answers \emph{what was being worked
under}; route~B answers \emph{what is this turn about}; route~C answers
\emph{where is this pattern written down}. As typed edges:

\begin{itemize}\itemsep0pt\parskip0pt
\item \textbf{A}: turn $\rightarrow$ clocked mission
\item \textbf{B}: turn $\rightarrow$ retrieved pattern
\item \textbf{C}: pattern $\rightarrow$ mentioning mission
\end{itemize}

\noindent A and B meet on the turn id, and that
coincidence supplies a derived association between a pattern and a mission.
It is not the authored citation used in the cascade of \S\ref{sec:factoring}. C reaches the same relation from
the other end, without needing the operator to have been clocked in correctly ---
or at all.

\begin{figure}[t]
  \centering
  \includegraphics[width=\linewidth]{turn-pattern-mission-chain}
  \caption{Four turns, each with its three answers. Route~B (solid, coloured)
  goes through the pattern; route~C (brown) continues from that pattern to the
  missions whose documents name it; route~A (dashed) goes straight from the turn
  to the mission the operator was clocked into. A and C reach the same kind of node through different relations. Green marks a pattern the curated layer already
  cites, blue one it never has. The following prose gives the generated full-window counts.}
  \label{fig:turn-chain}
\end{figure}

In the window measured for Figure~\ref{fig:turn-chain}, \ChainTurns{} operator
turns carried both annotations, yielding \ChainEdges{} turn$\rightarrow$pattern
edges over \ChainPatterns{} patterns and \ChainMissions{} missions. Of those
edges, \ChainAgree{} name a pattern the curated layer already cites; \ChainNew{}
do not.

Route C then reaches \ChainMentioned{} of those \ChainPatterns{} patterns,
adding \ChainMentionEdges{} pattern$\rightarrow$mission edges across
\ChainMentionMissions{} missions --- against the \ChainMissions{} that clock-in
alone supplies. And the two disagree completely: of the patterns route C
reaches, \ChainMentionAgree{} are mentioned in the mission the operator was
actually clocked into.

The routes answer different questions: clocked session context and textual
mention. Their disagreement does not identify an incorrect clock, cross-mission
work or an omitted citation without further evidence. Likewise, the
\ChainNew{} retrieved associations absent from the curated set of
\ChainCurated{} patterns identify candidates for inspection, not demonstrated
attention or missing warranted citations.

Three cautions, stated here because they bound what the figure supports.

First, retrieval is not use. Embedding proximity proposes a topical association; it does not establish
that the pattern was applied. An application receipt would need separate evidence; nothing in this
section supplies it.

Second, routes~A and~B overlap only where both fire. Route~A's coverage is
partial, so \ChainTurns{} is a floor set by the weaker mechanism, not a measure
of how much of the operator's attention is traceable in principle. Improving the
clock raises it directly --- and route~C sidesteps it, which is much of the point
of having a third route rather than a better second one.

Third, and most important for what may be built on this: these derived edges
must not be merged into the curated layer. Their provenance is different --- a
cosine score, not a human citation --- and merging them would destroy the
distinction that makes the gap measurable at all. They belong in a separate
relation carrying score and turn id, where agreement with a curated citation is
a cross-check rather than an assumption.

With those caveats the claim is modest and checkable: the operator is not
outside the record. He is in it, per turn, at a rate that can be counted, and
the chain from a turn he took to a mission he was working on runs through the
pattern library. \S\ref{sec:operator} argues what to do with that; this section
establishes only that there is something to do it with.
